\begin{proposition}[C13: \status{Proved}; standard tails are not a weakly initial family]\label{sa08:C13}
Work over $k=\Q$, and let $F$ be the standard theory. For $N\ge1$
let $A_N\subset F$ have coefficient algebra $S$ in weights below
$N$ and agree with $F$ in all weights at least $N$. These are
admissible theories, but for every $N$ there is no morphism
\[
 A_N\longrightarrow G_2,
\]
where $G_2=S\otimes_{\Q\{U\}}\mathrm{DR}(\Phi,U)$,
$\operatorname{wt}(U)=2$, and $U\mapsto T^{(2)}$.
The obstruction already holds for graded simplicial $S$-algebra
maps, without requiring compatibility with differentials or divided
powers. Thus replacing an arbitrary theory by a sufficiently high
standard tail cannot supply the generalized common-source theorem.
\end{proposition}
\begin{proof}
The definition of $A_N$ is closed under multiplication, the
differential and divided powers: a positive weight never decreases
under these operations, and its low-weight part is the coefficient
algebra. It resolves $S$ simplexwise, and its exterior-degree-$q$
term is flasque when the weight is at least $\max(N,q+1)$.
Thus it is admissible. Admissibility of $G_2$ follows from
Statement~\ref{sa08:C09}, since $\Q$ is a $\Z_{(2)}$-algebra.

In simplicial degree one and exterior degree zero, write
\[
 F_1^0=\Q[T,x],\qquad (G_2)_1^0=R=\Q[T,y],
 \qquad \operatorname{wt}(T)=\operatorname{wt}(x)=1,
 \quad\operatorname{wt}(y)=2.
\]
The two faces of $x$ are $0$ and $T$. Suppose a graded simplicial
$S$-algebra map $h:A_N\to G_2$ existed, and put
$a_n=h(x^n)\in R$ for $n\ge N$. One face of $a_n$ is $T^n$,
so every $a_n$ is nonzero. Multiplication gives
\[
 a_na_m=a_{n+m},\qquad
 a_{n+1}a_m=a_na_{m+1}\qquad(n,m\ge N).
\]
Consequently, in the fraction field of $R$, the ratio
$f=a_{n+1}/a_n$ is independent of $n$. We have
$a_{2N}=a_Nf^N=a_N^2$, whence $a_N=f^N$ and then $a_n=f^n$
for every $n\ge N$.

The element $f$ satisfies the monic equation $z^N-a_N=0$, so is
integral over $R$. The polynomial ring $R$ is integrally closed,
therefore $f\in R$. Since $f^N=a_N$ is homogeneous of weight $N$,
comparison of the least and greatest weights in the nonzero
polynomial $f$ shows that $f$ is homogeneous of weight one. Hence
$f=cT$ for some $c\in\Q$. Both simplicial faces of every $f^n$
are therefore the same. This contradicts the two faces $0$ and
$T^n$ required of $a_n=h(x^n)$. No such map exists.
\end{proof}
